\subsection{The chosen root in a connected family}
\label{sec:ftb-key-connected-root}

For a complex key datum in a normal form \((\tau,M,a,b,c)\), put
\[
                            t(x)=\frac{(\det M)^4}{c^8}.
 \tag{R1}\label{eq:ftb-key-root-value}
\]
We first justify that this is a value of the datum, not of the
chosen normal form.

\begin{srlemma}[Independence of the normal form]
\label{lem:ftb-key-normal-value-unique}
With the actual complex theta uniformization and the Poisson theta
transformation law, \eqref{eq:ftb-key-root-value} is independent of
the half-characteristic normal form.
\end{srlemma}
\begin{proof}
Two exponential uniformizations of the same complex abelian variety
differ by a complex linear map \(\phi\); their specified tangent
isomorphisms determine it.  It maps their period lattices to each
other.  The quotient of their local line frames glues, by the
identical transition functions, to a nowhere-zero entire \(e(z)\).
Both use the same origin rigidification, so \(e(0)=1\).
The theta identity has the form
\[
        c'\theta[a',b'](z,\tau')
                 =e(z)c\theta[a,b](\phi z,\tau).
 \tag{R2}\label{eq:ftb-normal-frame-compare}
\]
The integer lattice map is symplectic.  Here is the actual
automorphy check, which avoids assuming that a group isomorphism
preserves an unspecified polarization.  Write the images of its
\(2g\) period generators as \(m_s+\tau n_s\).
The two automorphy laws and \eqref{eq:ftb-normal-frame-compare}
give the difference equation for the nowhere-zero entire frame
quotient.  Taking a global holomorphic logarithm on \(\mathbb C^g\)
gives
\(q(z+p_s)-q(z)=\ell_s(z)+d_s\), with linear \(\ell_s\);
the possible \(2\pi i\) multiples are absorbed in \(d_s\).
Commuting two differences yields
\(\ell_s(p_t)=\ell_t(p_s)\).
Substitute the linear terms \(-2\pi i n_s\cdot\phi z\)
and their primed counterparts from \eqref{eq:ftb-level-aut}.
The quadratic \(\tau\) terms cancel by symmetry, giving equality
of the two integral alternating forms
\(n_s\cdot m_t-n_t\cdot m_s\).
Thus the integer lattice matrix is symplectic.

Consequently \(\tau'=\delta\tau\) and
\(\phi z=j^tz\), \(j=C\tau+D\), for that symplectic matrix.
The theta transformation gives the multiplier \(r\) with
\(r^8=(\det j)^4\).
After its Gaussian exponential is absorbed, compare the two
theta functions at the same period.  The entire quotient \(\varrho\) between these same-period
functions has sign multipliers on each unit and \(\tau\)-period,
because both characteristics are half-integral.
Thus \(\varrho^2\) has every period of the full lattice.
It is bounded on a compact real period parallelepiped, hence
bounded on \(\mathbb C^g\), and is constant by Liouville.
Since \(\varrho\) is nowhere zero, differentiating its constant
square shows that \(\varrho\) itself is constant.
All its sign multipliers are therefore \(1\).
The two characteristics agree modulo integers; our representatives
\(0,\tfrac12\) make them equal.  Its value at zero then gives
\(c^8=(c')^8(\det j)^4\).
The tangent comparison gives \(M=M'j\).
Taking determinants proves
\((\det M)^4/c^8=(\det M')^4/(c')^8\).
\end{proof}

We will use the following analytic family input, with its exact
scope.  On a small ball of parameters for a smooth proper polarized
abelian scheme, its relative exponential and the chosen
dual-number tangent frame admit jointly holomorphic local
coordinates; a fixed integral symplectic basis of the period
lattice continues holomorphically; its torsion sections and Weil
pairing continue with the same labels.  The resulting normalized
period has positive imaginary part, and the theta section has,
on each line chart, its actual local nowhere-zero quotient by
one of the finitely many half-characteristic theta functions.
These are respectively the source's relative exponential,
period-continuation, torsion/Weil, and normal-fibre statements.
We prove the finite-choice and invariant-continuity step below.
Construction of this family input is distinct from the fixed
period holomorphy in
Proposition~\ref{prop:ftb-key-holomorphic-period}.

\begin{srlemma}[The source's circle quotient argument]
\label{lem:ftb-circle-quotient}
Let \(f(b,z),\theta(b,z)\) be holomorphic on an open subset of
finite-dimensional complex parameter and fibre spaces.
Suppose \(\theta(b_0,-)\) is not the zero germ at \(z_0\).
If, for parameters in a set \(S\), the fibre function \(f/\theta\)
extends holomorphically, then near \((b_0,z_0)\) there is a
holomorphic \(G(b,z)\) satisfying \(G\theta=f\) on parameters
in \(\overline S\).
If a quotient \(F\theta=f\) is defined on all fibres and is
holomorphic in each fibre, then it is jointly holomorphic.
\end{srlemma}
\begin{proof}
Choose a complex direction \(v\) on which the germ of
\(\theta(b_0,z_0+tv)\) is not identically zero.
Such a direction can also be obtained by joining \(z_0\) to a
nearby point where \(\theta\ne0\).
Its zeros on that line are isolated, so choose a sufficiently
small \(r>0\) for which theta is nonzero on \(|t|=r\).
Compactness of the circle and continuity keep it nonzero on
that same circle for all \((b,z)\) near \((b_0,z_0)\), and
keep the whole closed disk inside the chart.
Set
\[
 G(b,z)=\frac1{2\pi}\int_0^{2\pi}
       \frac{f(b,z+re^{iu}v)}{\theta(b,z+re^{iu}v)}\,du .
 \tag{R3}\label{eq:ftb-circle-average}
\]
Its denominator is bounded away from zero locally on the
compact circle.  Differentiating under that compact integral
makes \(G\) jointly holomorphic.
For \(b\in S\), the holomorphic fibre quotient has its value at
\(z\) equal to this circle average, by the one-variable
holomorphic mean-value formula.  Thus \(G\theta=f\) there.
Continuity, applied at each fixed \(z\), extends this equality
to \(b\in\overline S\).
In the second assertion the same circle formula equals \(F\)
on every nearby fibre, proving joint holomorphy.
No extension by a generic-fibre value is used.
\end{proof}

\begin{srlemma}[Finite theta choices and the rigidified scalar]
\label{lem:ftb-finite-frame-continuity}
Assume the actual local family input just stated.
The half-characteristic occurring in its normalized theta frames
is locally constant, and the scalar \(c(b)\) fixed by the
origin rigidification is continuous.
Consequently \(t(x_b)\) is continuous.
\end{srlemma}
\begin{proof}
In each line chart let \(f_j(b,z)\) be the actual algebraic
theta section read on the relative exponential.
It is holomorphic in both variables.
For each fibre, its quotient by the selected theta function
is locally a holomorphic unit.  Nonzero theta is nonzero on
a dense subset of each fibre by the identity theorem, so
these local quotients glue uniquely.

Suppose a different member \(k\) of the finite
half-characteristic set occurs at parameters approaching
\(b_0\).  Apply Lemma~\ref{lem:ftb-circle-quotient} to
\(f_j/\theta_k\), and also to its inverse
\(\theta_k/f_j\), along these parameters.
At the limiting fibre the two holomorphic quotients multiply
to \(1\): this equality holds on the dense set where the
functions are nonzero and then everywhere by continuity.
Compare with the unit quotient for the characteristic selected
at \(b_0\).  On overlapping line charts these quotients glue,
since their transition functions agree; dense nonvanishing
makes the gluing unique even at theta zeros.
They give a nowhere-zero entire ratio between theta functions
of characteristics \(k\) and \(k_0\).
The characteristic uniqueness in the preceding normal-form
argument forces \(k=k_0\), a contradiction.
The fibres of this finite characteristic choice are therefore
closed locally, hence open locally, proving local constancy.

With that characteristic fixed,
Lemma~\ref{lem:ftb-circle-quotient} makes the chart quotient
jointly holomorphic.
At \(z=0\) its value is \(c(b)\) times the actual line
rigidification value \(r_j(b)\).
That value is holomorphic and nonzero.
Dividing gives continuity of \(c(b)\), including when the
theta function itself vanishes at the origin.

For completeness, the tangent matrix is continuous as well.
Let \(P(b)\) be the continued first half of the period basis;
its inverse \(h(b)\) is holomorphic by the adjugate formula.
The vector \(W_i(b)\) in exponential coordinates representing
the \(i\)-th dual-number tangent is holomorphic: choose
\(g\) local functions whose derivatives give an invertible
Jacobian at the identity, and solve that finite system.
Both its coefficients and its tangent evaluations are
holomorphic in \(b\).  The normal-form matrix has rows
\(2\pi i\,h(b)W_i(b)\); thus its determinant is continuous.
It is nonzero, and \(c(b)\ne0\).
Formula \eqref{eq:ftb-key-root-value} proves continuity of \(t\).
\end{proof}

\begin{srlemma}[A powered rule on a smooth base is an actual unit]
\label{lem:ftb-powered-rule-smooth-unit}
Let \(k\) have characteristic zero, let \(S\) be a smooth integral
finitely generated \(k\)-algebra, and let \(y/S\) be a key datum.
The powered rule of
Proposition~\ref{prop:ftb-normalized-powered-key} is represented
by a unit \(u\in S^\times\), at every field-valued point of \(S\).
In particular
\[
                         t(y_s)^{N(g)}=s(u)
 \tag{R4}\label{eq:ftb-family-powered-unit}
\]
at all complex points.
\end{srlemma}
\begin{proof}
At every height-one prime the normal ring \(S\) localizes to a DVR.
The generic rule value and its inverse belong to that DVR
by its unit law.  The height-one intersection theorem for a
Noetherian normal domain puts both in \(S\), so they are
mutual inverses there.  This gives \(u\).

We must still identify its specialization in higher codimension.
The source proves this over a regular local ring \(A\) by
induction on its dimension.
Dimension zero is a field.  Otherwise choose
\(\pi\in\mathfrak m_A\setminus\mathfrak m_A^2\).
The quotient \(A/(\pi)\) is a regular local domain of dimension
one less, and \(A_{(\pi)}\) is a DVR.
The rule's characteristic-zero DVR specialization identifies
the generic value on \(A/(\pi)\) with \(u\bmod\pi\).
Apply the induction hypothesis to that quotient and its closed
point.  Every field-valued point with kernel \(\mathfrak m_A\)
factors through its residue field, so field functoriality gives
the assertion at that point.
Localizing \(S\) at any prime reduces its point to this case.
Thus the unit represents the rule at all field-valued points,
not only generically.
The exact complex normal value gives \eqref{eq:ftb-family-powered-unit}.
\end{proof}

\begin{srlemma}[A root cover with two lifted points]
\label{lem:ftb-key-connected-root-cover}
Let \(K\) be algebraically closed of characteristic zero,
\(R/K\) a smooth integral finitely generated algebra, and
\(u\in R^\times\).  Let \(n\geq1\), \(p_0,p:R\to K\), and
\(a_0^n=p_0(u)\).
There is a smooth integral \(K\)-algebra \(B\), a map
\(\phi:R\to B\), a unit \(r\in B^\times\) with
\(r^n=\phi(u)\), and lifts \(q_0,q:B\to K\) of \(p_0,p\),
such that \(q_0(r)=a_0\).
\end{srlemma}
\begin{proof}
Use induction on the prime factors of \(n\).
For a prime \(l\), if \(u=\rho^l\) in \(R\), take \(B=R\)
and replace \(\rho\) by
\((a_0/p_0(\rho))\rho\); its scalar is an \(l\)-th root of
unity and fixes its \(l\)-th power.
Otherwise \(X^l-u\) is irreducible over \(\operatorname{Frac}R\).
Indeed that field contains all \(l\)-th roots of unity; the
conjugates of a root form an orbit of size \(1\) or \(l\).
In the first case the root lies in the fraction field, is
integral over \(R\), and hence lies in \(R\), contrary to the
case assumption.
Take \(B=R[X]/(X^l-u)\).  It is a domain.
Its root is a unit, so its derivative \(lX^{l-1}\) is a unit;
the algebra is finite etale over \(R\), hence smooth over \(K\).
Algebraic closedness gives a lift of \(p\), and the specified
root \(a_0\) gives the desired lift of \(p_0\).
Apply this step and then the inductive root construction for
the remaining prime factors.  The specified root at the first
point is retained at each step.  This is the source's connected
root-cover construction; an arbitrary disconnected root torsor
would not give the required integral base.
\end{proof}

\begin{srlemma}[The root is fixed by one point]
\label{lem:ftb-root-fixed-by-point}
Let \(B\) be a smooth integral finitely generated complex algebra,
\(y/B\) a key datum, \(r\in B^\times\), and \(n\geq1\).
Suppose the actual normal value satisfies
\(t(y_q)^n=q(r)^n\) at every complex point \(q\).
If \(t(y_{q_0})=q_0(r)\) at one point, then
\(t(y_q)=q(r)\) at every point.
\end{srlemma}
\begin{proof}
The complex points of an integral finite-type complex variety
are connected.  One standard proof, used in the source, takes
a Noether-normalization principal open with primitive-element
unramified description, proves that this open is connected
using irreducibility and analytic continuation off its
discriminant, and uses density of that principal open in
complex points.
Smooth algebraic charts are complex analytic charts.
Lemma~\ref{lem:ftb-finite-frame-continuity}, applied in these
charts, makes \(q\mapsto t(y_q)\) continuous.
Evaluation \(q\mapsto q(r)\) is continuous and nonzero.
Their ratio is therefore a continuous map to the finite set
of \(n\)-th roots of unity.  It is constant on the connected
space, and its value at \(q_0\) is \(1\).
\end{proof}

\paragraph{Exact arithmetic-family interface.}
We call a datum \(x/\mathbb C\) \emph{connected to a known point}
if there are a number field \(k\), a smooth geometrically
integral finitely generated \(k\)-algebra \(S\), a key datum
\(y/S\), maps \(s:S\to\mathbb C\) and \(s_0:S\to k\) with
\(s_0\) a \(k\)-map, and \(t_0\in k\), such that
\(x=s^*y\), and for every embedding \(e:k\to\mathbb C\),
\[
                         t((e s_0)^*y)=e(t_0).
 \tag{R5}\label{eq:ftb-connfam}
\]
All identifications here are key pullbacks: line bundle,
rigidification, theta section and ordered Lie frame are included.
Geometric integrality means \(F\otimes_k S\) is a domain for
\emph{every} field extension \(F/k\).
This is the literal source \(\mathrm{ConnFam}\), not just
path connectedness of an analytic moduli space.

\begin{srproposition}[Known families give equivariance under complex field maps]
\label{prop:ftb-key-complex-equivariance}
Assume the local analytic family input above, and that every
complex key datum is connected to a known point in the exact
sense \eqref{eq:ftb-connfam}.
Then for every field map \(f:\mathbb C\to\mathbb C\) and every
key pullback \(x'=f^*x\),
\[
                              t(x')=f(t(x)).
 \tag{R6}\label{eq:ftb-key-complex-funct}
\]
\end{srproposition}
\begin{proof}
For the family in \eqref{eq:ftb-connfam} put
\(e=s|_k\), and base change to
\(R=\mathbb C\otimes_{k,e}S\).
It is smooth integral; \(s\) and \(es_0\) give two complex
points \(p,p_0\).
Lemma~\ref{lem:ftb-powered-rule-smooth-unit} gives a unit
\(u\in S^\times\) with \(t^{N(g)}\) equal to its evaluations.
Apply Lemma~\ref{lem:ftb-key-connected-root-cover} to its
image in \(R\), choosing the root \(e(t_0)\) at \(p_0\).
It gives \(B,r,q,q_0\).
The root-fixed-by-one-point lemma makes \(t(x)=q(r)\).

Now repeat the entire base change by the possibly nonsurjective
embedding \(f:\mathbb C\to\mathbb C\).
The algebra \(\mathbb C\otimes_{\mathbb C,f}B\) is smooth and
integral: a finite-type integral algebra over an algebraically
closed field is geometrically integral.
The points \(fq,fq_0\) lift to this base change, and its root
is the image of \(r\).
At the lifted \(q_0\), \eqref{eq:ftb-connfam} gives the chosen
root \(f(e(t_0))\).
The same root-fixed-by-one-point lemma makes its value at the
lifted \(q\) equal to \(f(q(r))=f(t(x))\).
This is the normal value of \(x'\) by
Lemma~\ref{lem:ftb-key-normal-value-unique}.
No continuity of the field map \(f\) has been assumed or used.
\end{proof}

\paragraph{Frozen source concordance and scope.}
The base prefix is
\srcid{src/Development/AbelianVarietyTheory/Proofs/}.
Normal-value independence is
\srcid{ThetaLawProofs/ThetaFrameValueIndependentNormalForm.lean}
(\srcid{f823bc2749dfa176812b47f1be3902b284c5eada}).
The circle quotient and closure argument are
\srcid{ModuliProofs/DifferentiableOnQuotientOfDifferentiableOnSlice.lean}
(\srcid{efca854eca061af3cefd2f6ce3c9b9f13e6cc61c})
and \srcid{ModuliProofs/EqOfFiniteOfClosedFibres.lean}
(\srcid{7f7959eed9e909e48e267a075fa077ab74d989e9}).
The finite-character and origin-scalar argument is
\srcid{ThetaLawProofs/ExistsCoordinatesOfRelations.lean}
(\srcid{6da8aa085884d27e2caacc57952b26d8dd59704a}).
Its actual family assembly is
\srcid{ThetaLawProofs/ContinuousOnThetaFrameValueHolomorphicFamily.lean}
(\srcid{eae569828b1ee42499288b46d7fcf4dd39dc4062}):
the relative exponential, continued periods, and their
torsion/Weil compatibility remain the explicit family input
stated above.  These are produced in
\srcid{ModuliProofs/ExistsRelativeExpOfParts.lean},
\srcid{LatContMapContinuous.lean},
\srcid{TorsionEqAndWeilIsOfIsHolTo.lean}, and
\srcid{ThetaLawProofs/ExistsNormalOfWeil.lean}.
The smooth-base unit proof is
\srcid{ThetaLawProofs/ExistsUnitThetaKeyRuleSmoothBase.lean}
(\srcid{36413650a428826ff71462b79aded750faaa2d18}).
Connected root constancy is
\srcid{ThetaLawProofs/ThetaFrameValueEqUnitSmoothConnectedBase.lean}
(\srcid{aee5baa0b37f382a5ca6aeb4caca40e115b342aa}).
The integral root covers, geometric base change and equivariance
are \srcid{ThetaLawProofs/IsDomainTensorOfIsAlgClosed.lean}
(\srcid{df53eb7c2f461008fca381d1052984d14cc517d0}).
The arithmetic \(\mathrm{ConnFam}\) interface is expanded here
but its construction is a further producer; it is not inferred
from an arbitrary complex path.
